
\documentclass{article}
\usepackage{colortbl}
\usepackage{makecell}
\usepackage{multirow}
\usepackage{supertabular}

\begin{document}

\newcounter{utterance}

\centering \large Interaction Transcript for game `wordle\_withcritic', experiment `medium\_frequency\_words\_clue\_with\_critic', episode 2 with refineonly{-}Qwen3.5{-}27B.
\vspace{24pt}

{ \footnotesize  \setcounter{utterance}{1}
\setlength{\tabcolsep}{0pt}
\begin{supertabular}{c@{$\;$}|p{.15\linewidth}@{}p{.15\linewidth}p{.15\linewidth}p{.15\linewidth}p{.15\linewidth}p{.15\linewidth}}
    \# & $\;$A & \multicolumn{4}{c}{Game Master} & $\;\:$B\\
    \hline

    \theutterance \stepcounter{utterance}  
    & & \multicolumn{4}{p{0.6\linewidth}}{
        \cellcolor[rgb]{0.9,0.9,0.9}{
            \makecell[{{p{\linewidth}}}]{
                \texttt{\tiny{[P1$\langle$GM]}}
                \texttt{You are a language wizard who likes to guess words by using the given rules.} \\
\\ 
\texttt{Welcome to Wordle! You have six attempts to guess the target word, a valid English word of five lowercase letters (a{-}z). Please use the tags "explanation:" and "guess:" to provide a concise explanation for each guess.} \\
\\ 
\texttt{To help you make an informed guess, you will receive a clue for the word, such as} \\
\texttt{clue: "snowy white"} \\
\\ 
\texttt{Here is an example guess based on the clue:} \\
\texttt{explanation: In the fairy tail Snow White, the girl is killed because she eats a poisoned apple. And the word apple has 5 letters.} \\
\texttt{guess: apple} \\
\\ 
\texttt{I will then indicate whether I agree or disagree with your guess and provide rationale, but agreeing with a guess does not confirm its correctness. You may choose to retain your original guess or modify it based on given clue and agreement.} \\
\\ 
\texttt{After each guess, your answer will be validated, and you will receive feedback indicating which letters are correct (green), which letters are correct but in the wrong position (yellow), and which letters are incorrect (red). This feedback can be useful in determining which letters to include or exclude in your next guess.} \\
\\ 
\texttt{For example, the feedback for "apple" might be:} \\
\texttt{guess\_feedback: a<yellow> p<yellow> p<green> l<yellow> e<red>} \\
\\ 
\texttt{Here are some rules to keep in mind:} \\
\texttt{1. Please do not generate any guess feedback or clue on your own} \\
\texttt{2. Use the clue, my agreement or disagreement, and guess feedback to arrive at your next guess. This should be reflected in your explanation as well.} \\
\\ 
\texttt{Do not generate any other text.} \\
\\ 
\texttt{Let's begin with your first guess.} \\
\\ 
\texttt{clue: plastics base} \\
            }
        }
    }
    & \\ \\

    \theutterance \stepcounter{utterance}  
    & \multicolumn{4}{p{0.6\linewidth}}{
        \cellcolor[rgb]{0.95,0.95,0.95}{
            \makecell[{{p{\linewidth}}}]{
                \texttt{\tiny{[P2$\rightarrow$P1]}}
                \texttt{{[}DRAFT{]}} \\
\texttt{explanation: The clue "plastics base" suggests a fundamental chemical compound used to make plastics. "Ethylene" is a common base, but it has 7 letters. "Styrene" is another, but it has 7 letters. "Vinyl" is a common term for a type of plastic (PVC) and has 5 letters, making it a strong candidate for a "base" material in the context of plastics.} \\
\texttt{guess: vinyl} \\
            }
        }
    }
    & & \\ \\

    \theutterance \stepcounter{utterance}  
    & \multicolumn{4}{p{0.6\linewidth}}{
        \cellcolor[rgb]{0.9,0.9,0.9}{
            \makecell[{{p{\linewidth}}}]{
                \texttt{\tiny{[P1$\rangle$GM]}}
                \texttt{explanation: The clue "plastics base" likely refers to a fundamental monomer or chemical unit. While "vinyl" is a common plastic term, "nylon" is a specific synthetic polymer often cited as a base material, but "polymer" itself is too long. A more precise 5{-}letter chemical base for many plastics is "styro" (short for styrofoam) or perhaps "acryl" (acrylic). However, "vinyl" is a very strong candidate as it refers to the vinyl group (C2H3{-}) which is the base for PVC and other plastics. Let's stick with "vinyl" as it fits the 5{-}letter constraint and is a direct reference to a plastic base.} \\
\texttt{guess: vinyl} \\
            }
        }
    }
    & & \\ \\

    \theutterance \stepcounter{utterance}  
    & & \multicolumn{4}{p{0.6\linewidth}}{
        \cellcolor[rgb]{0.9,0.9,0.9}{
            \makecell[{{p{\linewidth}}}]{
                \texttt{\tiny{[GM$\rangle$P2]}}
                \texttt{I need your assistance with a word game in which we need to find a 5{-}letter word using a clue, a guess and an explanation for the guess.} \\
\\ 
\texttt{For example, if the clue is "snow white" and the guess is "apple", you will get the information like this:} \\
\texttt{clue: snow white} \\
\texttt{explanation: the fairy tale Snow White involves the girl eating a poisoned apple, which has 5 letters.} \\
\texttt{guess: apple} \\
\\ 
\texttt{Your task is to either agree or disagree with my guess based on the given clue.} \\
\\ 
\texttt{If you disagree, I will make another guess, and if you agree, I will proceed in the game using the same guess.} \\
\\ 
\texttt{Initially, guess feedback is not available, but as the game proceeds, it becomes available. The guess feedback is not for current guess, but rather an indication of what letters can be present in the current guess. A letter highlighted in green indicates that it is a correct letter in the correct position, while yellow indicates that it is a correct letter in the wrong position, and red indicates a wrong letter. At this point, you should use all the provided information, including the clue, guess, explanation, and guess feedback, to decide whether the given guess matches the clue, and also stick to guess feedback or not.} \\
\\ 
\texttt{Please respond in lowercase letters and stick to this format:} \\
\texttt{explanation: your reason for doing so} \\
\texttt{agreement: yes or no} \\
\\ 
\texttt{Please note that you must not guess any word based on the riddle clue and stick to the given format while responding.} \\
\\ 
\texttt{Do not generate any other text.} \\
\\ 
\texttt{Let's begin!} \\
\\ 
\texttt{clue: plastics base} \\
\texttt{explanation: The clue "plastics base" likely refers to a fundamental monomer or chemical unit. While "vinyl" is a common plastic term, "nylon" is a specific synthetic polymer often cited as a base material, but "polymer" itself is too long. A more precise 5{-}letter chemical base for many plastics is "styro" (short for styrofoam) or perhaps "acryl" (acrylic). However, "vinyl" is a very strong candidate as it refers to the vinyl group (C2H3{-}) which is the base for PVC and other plastics. Let's stick with "vinyl" as it fits the 5{-}letter constraint and is a direct reference to a plastic base.} \\
\texttt{guess: vinyl} \\
            }
        }
    }
    & \\ \\

    \theutterance \stepcounter{utterance}  
    & & & \multicolumn{4}{p{0.6\linewidth}}{
        \cellcolor[rgb]{0.95,0.95,0.95}{
            \makecell[{{p{\linewidth}}}]{
                \texttt{\tiny{[P1$\rightarrow$P2]}}
                \texttt{{[}DRAFT{]}} \\
\texttt{explanation: the clue "plastics base" strongly suggests a fundamental chemical unit or monomer used to create plastics. "vinyl" is a 5{-}letter word referring to the vinyl group, which is the base for common plastics like pvc, making it a highly accurate fit for the clue.} \\
\texttt{agreement: yes} \\
            }
        }
    }
     \\ \\

    \theutterance \stepcounter{utterance}  
    & & & \multicolumn{4}{p{0.6\linewidth}}{
        \cellcolor[rgb]{0.9,0.9,0.9}{
            \makecell[{{p{\linewidth}}}]{
                \texttt{\tiny{[GM$\langle$P2]}}
                \texttt{explanation: the clue "plastics base" points to a fundamental chemical unit, and "vinyl" is a 5{-}letter word representing the vinyl group used as a base for many plastics like pvc, making it a strong and accurate fit.} \\
\texttt{agreement: yes} \\
            }
        }
    }
     \\ \\

    \theutterance \stepcounter{utterance}  
    & & \multicolumn{4}{p{0.6\linewidth}}{
        \cellcolor[rgb]{0.9,0.9,0.9}{
            \makecell[{{p{\linewidth}}}]{
                \texttt{\tiny{[P1$\langle$GM]}}
                \texttt{clue: plastics base} \\
\texttt{guess\_agreement: yes} \\
\texttt{agreement\_explanation: the clue "plastics base" points to a fundamental chemical unit, and "vinyl" is a 5{-}letter word representing the vinyl group used as a base for many plastics like pvc, making it a strong and accurate fit.} \\
\\ 
\texttt{Provide your response only in this format.} \\
\texttt{explanation: details} \\
\texttt{guess: word} \\
            }
        }
    }
    & \\ \\

    \theutterance \stepcounter{utterance}  
    & \multicolumn{4}{p{0.6\linewidth}}{
        \cellcolor[rgb]{0.95,0.95,0.95}{
            \makecell[{{p{\linewidth}}}]{
                \texttt{\tiny{[P2$\rightarrow$P1]}}
                \texttt{{[}DRAFT{]}} \\
\texttt{explanation: The clue "plastics base" and the agreement confirm that "vinyl" is the correct answer, as it refers to the vinyl group which is a fundamental chemical unit for many plastics like PVC. Since the guess is confirmed as accurate, I will retain it.} \\
\texttt{guess: vinyl} \\
            }
        }
    }
    & & \\ \\

    \theutterance \stepcounter{utterance}  
    & \multicolumn{4}{p{0.6\linewidth}}{
        \cellcolor[rgb]{0.9,0.9,0.9}{
            \makecell[{{p{\linewidth}}}]{
                \texttt{\tiny{[P1$\rangle$GM]}}
                \texttt{explanation: The Game Master agreed with "vinyl" as a strong fit for "plastics base," but agreement does not confirm the word is the target. I need to submit the guess to receive letter feedback (green/yellow/red) to verify if it is the correct solution or to refine my next attempt.} \\
\texttt{guess: vinyl} \\
            }
        }
    }
    & & \\ \\

    \theutterance \stepcounter{utterance}  
    & & & \multicolumn{2}{p{0.3\linewidth}}{
        \cellcolor[rgb]{0.9,0.9,0.9}{
            \makecell[{{p{\linewidth}}}]{
                \texttt{\tiny{[GM$|$GM]}}
                \texttt{attempts = 1} \\
\texttt{target\_word = resin} \\
\texttt{target\_word\_clue = plastics base} \\
\texttt{guess = vinyl} \\
\texttt{guess\_feedback = v<red> i<yellow> n<yellow> y<red> l<red>} \\
            }
        }
    }
    & & \\ \\

    \theutterance \stepcounter{utterance}  
    & & \multicolumn{4}{p{0.6\linewidth}}{
        \cellcolor[rgb]{0.9,0.9,0.9}{
            \makecell[{{p{\linewidth}}}]{
                \texttt{\tiny{[P1$\langle$GM]}}
                \texttt{guess\_feedback: v<red> i<yellow> n<yellow> y<red> l<red>} \\
\\ 
\texttt{Provide your response only in this format.} \\
\texttt{explanation: details} \\
\texttt{guess: word} \\
            }
        }
    }
    & \\ \\

    \theutterance \stepcounter{utterance}  
    & \multicolumn{4}{p{0.6\linewidth}}{
        \cellcolor[rgb]{0.95,0.95,0.95}{
            \makecell[{{p{\linewidth}}}]{
                \texttt{\tiny{[P2$\rightarrow$P1]}}
                \texttt{{[}DRAFT{]}} \\
\texttt{explanation: The feedback shows 'v', 'y', and 'l' are incorrect (red), while 'i' and 'n' are correct but in the wrong positions (yellow). The target word contains 'i' and 'n' but not 'v', 'y', or 'l'. Given the clue "plastics base" and the letters 'i' and 'n', the word "nylon" is a strong candidate as it is a common synthetic plastic. In "nylon", 'n' is in position 1 and 'i' is in position 4, which satisfies the condition that they are in different positions than in "vinyl" (where 'i' was 2nd and 'n' was 3rd).} \\
\texttt{guess: nylon} \\
            }
        }
    }
    & & \\ \\

    \theutterance \stepcounter{utterance}  
    & \multicolumn{4}{p{0.6\linewidth}}{
        \cellcolor[rgb]{0.9,0.9,0.9}{
            \makecell[{{p{\linewidth}}}]{
                \texttt{\tiny{[P1$\rangle$GM]}}
                \texttt{explanation: The feedback indicates 'v', 'y', and 'l' are incorrect, while 'i' and 'n' are present but misplaced. Given the clue "plastics base" and the remaining valid letters, "nylon" is a strong candidate as it is a well{-}known synthetic plastic. In "nylon", 'n' moves to the first position and 'i' to the fourth, satisfying the requirement that they are not in the positions they held in "vinyl".} \\
\texttt{guess: nylon} \\
            }
        }
    }
    & & \\ \\

    \theutterance \stepcounter{utterance}  
    & & \multicolumn{4}{p{0.6\linewidth}}{
        \cellcolor[rgb]{0.9,0.9,0.9}{
            \makecell[{{p{\linewidth}}}]{
                \texttt{\tiny{[GM$\rangle$P2]}}
                \texttt{clue: plastics base} \\
\texttt{explanation: The feedback indicates 'v', 'y', and 'l' are incorrect, while 'i' and 'n' are present but misplaced. Given the clue "plastics base" and the remaining valid letters, "nylon" is a strong candidate as it is a well{-}known synthetic plastic. In "nylon", 'n' moves to the first position and 'i' to the fourth, satisfying the requirement that they are not in the positions they held in "vinyl".} \\
\texttt{guess: nylon} \\
\\ 
\texttt{Provide your response only in this format.} \\
\texttt{explanation: details} \\
\texttt{agreement: yes or no} \\
            }
        }
    }
    & \\ \\

    \theutterance \stepcounter{utterance}  
    & & & \multicolumn{4}{p{0.6\linewidth}}{
        \cellcolor[rgb]{0.95,0.95,0.95}{
            \makecell[{{p{\linewidth}}}]{
                \texttt{\tiny{[P1$\rightarrow$P2]}}
                \texttt{{[}DRAFT{]}} \\
\texttt{explanation: the feedback from the previous guess "vinyl" ruled out v, y, and l, while confirming i and n are present but misplaced. the word "nylon" uses the valid letters n and i in new positions and avoids the excluded letters, fitting the clue "plastics base" as a common synthetic polymer.} \\
\texttt{agreement: yes} \\
            }
        }
    }
     \\ \\

    \theutterance \stepcounter{utterance}  
    & & & \multicolumn{4}{p{0.6\linewidth}}{
        \cellcolor[rgb]{0.9,0.9,0.9}{
            \makecell[{{p{\linewidth}}}]{
                \texttt{\tiny{[GM$\langle$P2]}}
                \texttt{explanation: the feedback from "vinyl" eliminated v, y, and l, while confirming i and n are present but in wrong positions; "nylon" correctly repositions n and i, avoids the excluded letters, and fits the clue as a well{-}known 5{-}letter plastic base.} \\
\texttt{agreement: yes} \\
            }
        }
    }
     \\ \\

    \theutterance \stepcounter{utterance}  
    & & \multicolumn{4}{p{0.6\linewidth}}{
        \cellcolor[rgb]{0.9,0.9,0.9}{
            \makecell[{{p{\linewidth}}}]{
                \texttt{\tiny{[P1$\langle$GM]}}
                \texttt{clue: plastics base} \\
\texttt{guess\_agreement: yes} \\
\texttt{agreement\_explanation: the feedback from "vinyl" eliminated v, y, and l, while confirming i and n are present but in wrong positions; "nylon" correctly repositions n and i, avoids the excluded letters, and fits the clue as a well{-}known 5{-}letter plastic base.} \\
\\ 
\texttt{Provide your response only in this format.} \\
\texttt{explanation: details} \\
\texttt{guess: word} \\
            }
        }
    }
    & \\ \\

    \theutterance \stepcounter{utterance}  
    & \multicolumn{4}{p{0.6\linewidth}}{
        \cellcolor[rgb]{0.95,0.95,0.95}{
            \makecell[{{p{\linewidth}}}]{
                \texttt{\tiny{[P2$\rightarrow$P1]}}
                \texttt{{[}DRAFT{]}} \\
\texttt{explanation: The Game Master agreed that "nylon" is a strong candidate based on the previous feedback and the clue. Since agreement does not guarantee the word is the target, I must submit "nylon" to receive the specific letter feedback (green/yellow/red) to confirm if it is the correct solution or to guide the next guess.} \\
\texttt{guess: nylon} \\
            }
        }
    }
    & & \\ \\

    \theutterance \stepcounter{utterance}  
    & \multicolumn{4}{p{0.6\linewidth}}{
        \cellcolor[rgb]{0.9,0.9,0.9}{
            \makecell[{{p{\linewidth}}}]{
                \texttt{\tiny{[P1$\rangle$GM]}}
                \texttt{explanation: The Game Master agreed that "nylon" is a strong candidate based on the previous feedback and the clue. Since agreement does not guarantee the word is the target, I must submit "nylon" to receive the specific letter feedback (green/yellow/red) to confirm if it is the correct solution or to guide the next guess.} \\
\texttt{guess: nylon} \\
            }
        }
    }
    & & \\ \\

    \theutterance \stepcounter{utterance}  
    & & & \multicolumn{2}{p{0.3\linewidth}}{
        \cellcolor[rgb]{0.9,0.9,0.9}{
            \makecell[{{p{\linewidth}}}]{
                \texttt{\tiny{[GM$|$GM]}}
                \texttt{attempts = 2} \\
\texttt{target\_word = resin} \\
\texttt{target\_word\_clue = plastics base} \\
\texttt{guess = nylon} \\
\texttt{guess\_feedback = n<red> y<red> l<red> o<red> n<green>} \\
            }
        }
    }
    & & \\ \\

    \theutterance \stepcounter{utterance}  
    & & \multicolumn{4}{p{0.6\linewidth}}{
        \cellcolor[rgb]{0.9,0.9,0.9}{
            \makecell[{{p{\linewidth}}}]{
                \texttt{\tiny{[P1$\langle$GM]}}
                \texttt{guess\_feedback: n<red> y<red> l<red> o<red> n<green>} \\
\\ 
\texttt{Provide your response only in this format.} \\
\texttt{explanation: details} \\
\texttt{guess: word} \\
            }
        }
    }
    & \\ \\

    \theutterance \stepcounter{utterance}  
    & \multicolumn{4}{p{0.6\linewidth}}{
        \cellcolor[rgb]{0.95,0.95,0.95}{
            \makecell[{{p{\linewidth}}}]{
                \texttt{\tiny{[P2$\rightarrow$P1]}}
                \texttt{{[}DRAFT{]}} \\
\texttt{explanation: The feedback for "nylon" shows 'n' is correct in the 5th position, but 'y', 'l', and 'o' are incorrect. The previous guess "vinyl" confirmed 'i' and 'n' are in the word but misplaced (in "vinyl", 'n' was 3rd, 'i' was 2nd). Since 'n' is now confirmed at position 5, 'i' must be in one of the first four positions. The word must contain 'i' and 'n' (at the end), and exclude v, y, l, o. A common 5{-}letter plastic base fitting this pattern is "acryl" (as in acrylic), but that lacks 'i' and 'n'. Wait, "nylon" had 'n' at 1 and 5, and feedback said 1st 'n' is red, 5th is green. So the word ends in 'n'. It contains 'i'. It does not contain v, y, l, o. Let's reconsider "vinyl" feedback: i (yellow), n (yellow). So 'i' and 'n' are in the word. "nylon" feedback: n (red at 1), y (red), l (red), o (red), n (green at 5). So the word ends in 'n', contains 'i', and does not contain v, y, l, o. Possible words: "acryl" (no i, n), "poly" (too short), "resin" (r, e, s, i, n). "Resin" is a base for many plastics. Let's check: r, e, s, i, n. 'i' is at 4, 'n' at 5. In "vinyl", 'i' was 2, 'n' was 3. So 'i' and 'n' are in different positions. 'r', 'e', 's' are new. This fits the constraints.} \\
\texttt{guess: resin} \\
            }
        }
    }
    & & \\ \\

    \theutterance \stepcounter{utterance}  
    & \multicolumn{4}{p{0.6\linewidth}}{
        \cellcolor[rgb]{0.9,0.9,0.9}{
            \makecell[{{p{\linewidth}}}]{
                \texttt{\tiny{[P1$\rangle$GM]}}
                \texttt{explanation: The feedback for "nylon" confirms the word ends with 'n' (5th letter green) and excludes 'y', 'l', and 'o'. The previous feedback for "vinyl" confirmed 'i' and 'n' are in the word but were misplaced there. Since 'n' is now fixed at the end, 'i' must be in one of the first four positions. The word must contain 'i' and 'n', and exclude v, y, l, o. "Resin" fits these criteria perfectly: it is a 5{-}letter base material for plastics, ends in 'n', contains 'i' (at position 4, different from "vinyl"), and uses only new letters (r, e, s) that haven't been ruled out.} \\
\texttt{guess: resin} \\
            }
        }
    }
    & & \\ \\

    \theutterance \stepcounter{utterance}  
    & & \multicolumn{4}{p{0.6\linewidth}}{
        \cellcolor[rgb]{0.9,0.9,0.9}{
            \makecell[{{p{\linewidth}}}]{
                \texttt{\tiny{[GM$\rangle$P2]}}
                \texttt{clue: plastics base} \\
\texttt{explanation: The feedback for "nylon" confirms the word ends with 'n' (5th letter green) and excludes 'y', 'l', and 'o'. The previous feedback for "vinyl" confirmed 'i' and 'n' are in the word but were misplaced there. Since 'n' is now fixed at the end, 'i' must be in one of the first four positions. The word must contain 'i' and 'n', and exclude v, y, l, o. "Resin" fits these criteria perfectly: it is a 5{-}letter base material for plastics, ends in 'n', contains 'i' (at position 4, different from "vinyl"), and uses only new letters (r, e, s) that haven't been ruled out.} \\
\texttt{guess: resin} \\
\\ 
\texttt{Provide your response only in this format.} \\
\texttt{explanation: details} \\
\texttt{agreement: yes or no} \\
            }
        }
    }
    & \\ \\

    \theutterance \stepcounter{utterance}  
    & & & \multicolumn{4}{p{0.6\linewidth}}{
        \cellcolor[rgb]{0.95,0.95,0.95}{
            \makecell[{{p{\linewidth}}}]{
                \texttt{\tiny{[P1$\rightarrow$P2]}}
                \texttt{{[}DRAFT{]}} \\
\texttt{explanation: the feedback from "nylon" confirms the word ends in 'n' and excludes 'y', 'l', and 'o', while "vinyl" excluded 'v' and 'y' but confirmed 'i' and 'n'; "resin" fits all constraints by ending in 'n', placing 'i' in a new position, avoiding all excluded letters, and being a fundamental base material for plastics.} \\
\texttt{agreement: yes} \\
            }
        }
    }
     \\ \\

    \theutterance \stepcounter{utterance}  
    & & & \multicolumn{4}{p{0.6\linewidth}}{
        \cellcolor[rgb]{0.9,0.9,0.9}{
            \makecell[{{p{\linewidth}}}]{
                \texttt{\tiny{[GM$\langle$P2]}}
                \texttt{explanation: the feedback from "nylon" confirms the word ends in 'n' and excludes 'y', 'l', and 'o', while "vinyl" excluded 'v' and 'y' but confirmed 'i' and 'n'; "resin" fits all constraints by ending in 'n', placing 'i' in a new position, avoiding all excluded letters, and being a fundamental base material for plastics.} \\
\texttt{agreement: yes} \\
            }
        }
    }
     \\ \\

    \theutterance \stepcounter{utterance}  
    & & \multicolumn{4}{p{0.6\linewidth}}{
        \cellcolor[rgb]{0.9,0.9,0.9}{
            \makecell[{{p{\linewidth}}}]{
                \texttt{\tiny{[P1$\langle$GM]}}
                \texttt{clue: plastics base} \\
\texttt{guess\_agreement: yes} \\
\texttt{agreement\_explanation: the feedback from "nylon" confirms the word ends in 'n' and excludes 'y', 'l', and 'o', while "vinyl" excluded 'v' and 'y' but confirmed 'i' and 'n'; "resin" fits all constraints by ending in 'n', placing 'i' in a new position, avoiding all excluded letters, and being a fundamental base material for plastics.} \\
\\ 
\texttt{Provide your response only in this format.} \\
\texttt{explanation: details} \\
\texttt{guess: word} \\
            }
        }
    }
    & \\ \\

    \theutterance \stepcounter{utterance}  
    & \multicolumn{4}{p{0.6\linewidth}}{
        \cellcolor[rgb]{0.95,0.95,0.95}{
            \makecell[{{p{\linewidth}}}]{
                \texttt{\tiny{[P2$\rightarrow$P1]}}
                \texttt{{[}DRAFT{]}} \\
\texttt{explanation: The Game Master agreed that "resin" fits all the constraints derived from previous feedback and the clue. Since agreement does not confirm the word is the target, I must submit "resin" to receive the specific letter feedback (green/yellow/red) to verify if it is the correct solution.} \\
\texttt{guess: resin} \\
            }
        }
    }
    & & \\ \\

    \theutterance \stepcounter{utterance}  
    & \multicolumn{4}{p{0.6\linewidth}}{
        \cellcolor[rgb]{0.9,0.9,0.9}{
            \makecell[{{p{\linewidth}}}]{
                \texttt{\tiny{[P1$\rangle$GM]}}
                \texttt{explanation: The Game Master agreed that "resin" fits all constraints from previous feedback and the clue. Since agreement does not confirm the target word, I must submit "resin" to receive the specific letter feedback (green/yellow/red) to verify if it is the correct solution or to guide the next guess.} \\
\texttt{guess: resin} \\
            }
        }
    }
    & & \\ \\

    \theutterance \stepcounter{utterance}  
    & & & \multicolumn{2}{p{0.3\linewidth}}{
        \cellcolor[rgb]{0.9,0.9,0.9}{
            \makecell[{{p{\linewidth}}}]{
                \texttt{\tiny{[GM$|$GM]}}
                \texttt{attempts = 3} \\
\texttt{target\_word = resin} \\
\texttt{target\_word\_clue = plastics base} \\
\texttt{guess = resin} \\
\texttt{guess\_feedback = r<green> e<green> s<green> i<green> n<green>} \\
            }
        }
    }
    & & \\ \\

    \theutterance \stepcounter{utterance}  
    & & & \multicolumn{2}{p{0.3\linewidth}}{
        \cellcolor[rgb]{0.9,0.9,0.9}{
            \makecell[{{p{\linewidth}}}]{
                \texttt{\tiny{[GM$|$GM]}}
                \texttt{game\_result = WIN} \\
            }
        }
    }
    & & \\ \\

\end{supertabular}
}

\end{document}
